\documentclass{article}
\usepackage[utf8]{inputenc}

\usepackage[a4paper, hmargin=3cm, top=3cm, bottom=3cm]{geometry}

%\usepackage{euler} %usa la fuente "euler" para las ecuaciones
\usepackage{amsfonts} %para Z estilizada
\usepackage{pifont} %para más estilos de viñetas
\usepackage{amsmath} %necesario para equation*
\usepackage{amsthm} %necesario para proof enviroment
\usepackage{dsfont} %alguna letra caligráfica
\usepackage{tipa} %Epsilon bonito
\usepackage{graphicx} %para manejar imágenes
\usepackage{wrapfig} %para imágenes en modo "wrap"
\usepackage{lipsum} %para usar el texto de relleno
\usepackage{tikz-cd} %para los diagramas conmutativos
\usepackage{amssymb}
\usepackage{tabularx}
\graphicspath{ {./Images/} } %la dirección de la cual se sacan las imágenes
\usepackage{euler}
\usepackage[normalem]{ulem} % Subrayado
\usepackage{setspace} %Para interlineado
\usepackage{float}

\setlength{\parindent}{0pt}

\renewcommand{\figurename}{Figura}
\renewcommand{\labelitemii}{$\circ$}      % círculo en segundo nivel de itemize

\title{
  \vspace{-1.5cm}
  \textsc{
    \LARGE{Trabajo de Grado\\ Programa de 16 semanas}\\ \vspace{1cm}
    \large{Año 2026\\ \vspace{1cm} {Realizado por:} \\ \vspace{0.3cm} {Juan Camilo Lozano Suárez}}\\
  }
}
\date{}

\begin{document}
    
    \maketitle
    Esta propuesta se centra en la realización geométrica de los grafos, su integración en superficies y el estudio de sus simetrías mediante la teoría de grupos.
    \begin{itemize}
       \item \textbf{Fase 1: La topología del grafo (Semanas 1-4)} El bojetivo es transformar el objeto discreto (grafo) en un objeto continuo (espacio topológico) y estudiar su contenedor (superficie).
          \begin{itemize}
             \item \textbf{Semana 1: Realización geométrica.} Se define el grafo $G$ como un espacio topológico donde los vértices son puntos distintos y las aristas son subespacios homeomorfos al intervalo $[0,1]$.
             \item \textbf{Semana 2: El concepto de embebido.} Estudio del embebido como un homeomorfismo entre la realización geométrica de $G$ y un subespacio de $X$. Se analiza la distinción entre grafos planares y aquellos que requieren superficies de mayor complejidad.
             \item \textbf{Semana 3: Topología de superficies.} Clasificación de superficies (espacios compactos de Hausdorff localmente planares) mediante el uso de asas ($S_n$) y casquetes cruzados ($\tilde{S}_m$). Se introduce el concepto de orientabilidad.
             \item \textbf{Semana 4: Embebidos celulares y 2-celdas.} Definición de un embebido celular, donde cada componente de $X-G$ (cada cara) es una 2-celda abierta (disco abierto). Se estudia la importancia de que la cara sea topológicamente simple.
          \end{itemize}
       \item \textbf{Fase 2: Invariantes y dualidad (Semanas 5-8)} Se introducen herramientas matemáticas para clasificar los embebidos y sus estructuras recíprocas
          \begin{itemize}
             \item \textbf{Semana 5: La característica de Euler.} Estudio del invariante topológico $\chi(G=\# V- \# E + \# F)$. Se analiza cómo este valor depende exclusivamente de la superficie ($2-2n$ para orientables; $2-2m$ para no orientables).
             \item \textbf{Semana 6: Girth y largo de caras.} Relación entre la longitud del ciclo más corto (girth) y el perímetro de las caras. Uso de desigualdades ($\text{girth}\cdot \# F \leq 2 \# E$) para probar la imposibilidad de ciertos embebidos.
             \item \textbf{Semana 7: Dualidad geométrica.} Construcción del grafo dual colocando vértices en el interior de las 2-celdas del primal. Se estudia la propiedad de que el dual del dual recupera el grafo original.
             \item \textbf{Semana 8: Dualidad algebraica y rotaciones.} Introducción a los duales algebraicos (Whitney) basados en espacios de ciclos y cociclos. Se presenta el esquema de rotación (permutación cíclica de aristas) como la base algebraica para definir simetrías.
          \end{itemize}
       \item \textbf{Fase 3: Automorfismos y simetrías de mapas (Semanas 9-12)} Estudio de las funciones que preservan la estructura topológica del grafo embebido.
          \begin{itemize}
             \item \textbf{Semana 9: El grupo de automorfismos del mapa.} Definición de automorfismos que no solo preservan la adyacencia del grafo, sino también la rotación local en la superficie.
             \item \textbf{Semana 10: Mapas regulares.} Estudio de mapas con simetría excepcional donde el grupo de automorfismos actúa transitivamente sobre los arcos. Su orden es $2 \# E$ en casos de orientables.
             \item \textbf{Semana 11: Reflexividad y quirilidad.} Análisis de mapas reflexibles (isomorfos a su imagen especular mediante una simetría de espejo) y mapas quirales que solo admiten rotaciones.
             \item \textbf{Semana 12: Autodualidad e isomorfismos superiores.} Estudio de mapas autoduales, donde el mapa es isomorfo a su dual geométrico. Se analizan los mapas de Paley como ejemplos de alta simetría (regulares y autoduales).
          \end{itemize}
       \item \textbf{Fase 4: Estructuras algebraicas y topología geométrica (Semanas 13-16)} Conexión de los grafos con teoremas profundos de la topología de superficies.
          \begin{itemize}
             \item \textbf{Semana 13: El teorema de Hurwitz.} Estudio del límite algebraico para grupos de simetría en superficies: el orden máximo es $168(g-1)$ para superficies de género $g\geq 2$.
             \item \textbf{Semana 14: El teorema de Tucker.} Demostración de que cualquier grupo finito de homeomorfismos de una superficie puede representarse mediante la acción de simetría sobre un grafo de Cayley embebido.
             \item \textbf{Semana 15: Grafos de Cayley y voltajes.} Construcción de mapas regulares utilizando grupos y sus conjuntos generadores. Se introduce la teoría de grafos de voltaje para describir recubrimientos de superficies de forma económica.
             \item \textbf{Semana 16: Síntesis final.} Clasificación de mapas regulares en superficies de bajo género y consolidación de la relación entre la estructura del grupo de automorfismos y la topología de la superficie.
          \end{itemize}
    \end{itemize}

\end{document}
